% =====================================================================
%  Chapter 32 — The Record
% =====================================================================

\chapter{The Record}\idx{Record}
\label{ch:record}

This chapter reports the system's record and puts it beside four
comparators, with every number computed by the same code on the same
months. It is the chapter in which the reader decides whether the
architecture is worth running, and the decision is not the one the
book's author would make.

\section{The window}

The system runs from \numRecordStart{} to \numRecordEnd{}: \numRecordMonths{}
months, or just over twenty-four years. The panel it draws on is longer,
running from \numPanelStart{} to \numPanelEnd{}, because the model needs a
warm-up before it will produce a forecast and because the earlier years
are needed to establish the benchmark\idx{Benchmark}'s own history.

The initial capital is \numInitialCapital{} dollars, and the account is
rebalanced monthly.

\section{The five portfolios}

\begin{center}
\begin{tabular}{@{}P{1.62in}rrrrr@{}}
\toprule
 & \multicolumn{2}{c}{Growth} & & & \\
Variant & annual rate & Sharpe & drawdown & worst 12m & recovery\idx{Recovery} \\
\midrule
The index, bought and held   & \numSpyCAGR{} & \numSpySharpe{} & \numSpyDD{} & \numSpyWorstYear{} & \numSpyRecovery{} \\
Equal weight of every fund   & \numEqualCAGR{} & \numEqualSharpe{} & \numEqualDD{} & \numEqualWorstYear{} & \numEqualRecovery{} \\
The policy, no rebalancing\idx{Rebalancing}   & \numPolicyCAGR{} & \numPolicySharpe{} & \numPolicyDD{} & \numPolicyWorstYear{} & \numPolicyRecovery{} \\
\textbf{The system}           & \textbf{\numFullCAGR{}} & \textbf{\numFullSharpe{}} & \textbf{\numFullDD{}} & \textbf{\numFullWorstYear{}} & \numFullRecovery{} \\
The system, in a taxable account & \numTaxCAGR{} & \numTaxSharpe{} & \numTaxDD{} & \numTaxWorstYear{} & \numTaxRecovery{} \\
\bottomrule
\end{tabular}
\end{center}

Four things are visible in that table, and they pull in different
directions.

\textbf{The system does not make more money than the index.} It makes
\numFullCAGR{} against the index's \numSpyCAGR{}, a difference of about
two and a half percentage points a year. Over twenty-four years, a
difference of that size is the difference between the system finishing
with a little over six times the capital and the index finishing with
nearly nine.

\textbf{The system does not take more risk-adjusted.} Its Sharpe ratio is
\numFullSharpe{} against the index's \numSpySharpe{}, and its Sortino
ratio is \numFullSortino{} against \numSpySortino{}. On the
discipline-adjusted measure it is slightly ahead; on the risk-adjusted
measure it is a tie.

\textbf{The system does take less risk.} Its deepest drawdown is
\numFullDD{} against the index's \numSpyDD{}, its worst rolling
twelve-month outcome is \numFullWorstYear{} against the index's \numSpyWorstYear{}, and
its annualised volatility is \numFullVol{} against the index's
\numSpyVol{}. It also takes longer to recover --- \numFullRecovery{} months
against the index's \numSpyRecovery{}, which is an improvement, but the
taxable account's \numTaxRecovery{} months is not.

\textbf{An equal-weighted portfolio of all sixty-one funds does better on
the money and worse on the drawdown.} The naive comparator\idx{Comparator} earns
\numEqualCAGR{} with a drawdown of \numEqualDD{}. This is Chapter 25's
finding made concrete: the book's architecture is not the best available
way to invest in this universe, and a rule requiring no model, no
shrinkage, no Kelly solve and no governor earns more compound return over
this record.

\begin{center}\emph{The system's growth against the index's is shown in \figref{fig:wealth}, and the same two return streams as drawdowns in \figref{fig:drawdown}.}\end{center}

\section{The comparison the book does not make}

It would be easy, and wrong, to present the table above as though the
system were competitive. It is not competitive on compound return, and
the difference is large enough that a reader who wanted the most money
would not run this system.

What the system is competitive on is the pair that the owner's question
actually asks for: growth \emph{and} protection. Judged on the compound
growth\idx{Compound growth} rate alone, the index wins and the system loses. Judged on the
Sharpe ratio, the system and the index are indistinguishable. Judged on
the deepest drawdown, the system wins by a wide margin and the equal
weight comparator is second.

That is the trade this book offers, and it is a trade rather than a
dominance. A reader who wants maximum compound return should buy the
index and stop reading. A reader who wants the index's return with a
materially shallower worst episode should run the system, and should
understand that they are paying about two and a half percentage points a
year for it.

\begin{takeaway}[What the record says]
On growth the index wins by about two and a half percentage points a year.
On risk-adjusted return the two are indistinguishable. On the deepest
drawdown the system wins by seventeen percentage points. On worst
twelve-month outcome the system wins by fourteen. The book does not claim
the system is better; it claims it is different in one specific way, and
that the way may be worth paying for.
\end{takeaway}

\section{What the record does not cover}

Four exclusions, carried forward from the chapters that established them.

It does not cover the possibility that the record's particular period was
unlucky for the index or lucky for the system; Chapter 34 treats this
directly by recomputing the attribution from every starting month, and
the governor's conclusion survives it.

It does not cover a reader whose circumstances are different. The
alternative policies are in Chapter 33's family, and the defensive policy
earns \numDefCAGR{} with a drawdown of \numDefDD{} while the aggressive one
earns \numAggCAGR{} with a drawdown of \numAggDD{}.

It does not cover the costs of running the system. It does not cover the
behaviour of an investor who overrides it.

And it does not cover the possibility that the benchmark's twenty-four
years are the wrong benchmark for the reader's own horizon.

\begin{keytakeaways}
\item Over twenty-four years the system earns \numFullCAGR{} against the
index's \numSpyCAGR{}, with a Sharpe ratio of \numFullSharpe{} against
\numSpySharpe{} and a deepest drawdown of \numFullDD{} against \numSpyDD{}.
\item An equal-weighted portfolio of all sixty-one funds earns
\numEqualCAGR{} with a drawdown of \numEqualDD{} --- more money, more
pain, and no machinery.
\item The offer is a trade, not a dominance: the index's growth against
the system's shallower worst episode.
\item The record covers one period, one set of policies and one set of
costs; Chapter 34 recomputes the conclusion from every starting month.
\end{keytakeaways}
